\chapterbackmatter{Supplementary Exercises}

\begin{enumerate}
  \SupplementaryExercise{1}{Topology, Defect Codimension, and Causal Formation}
  {Inspired by T. W. B. Kibble, Topology of Cosmic Domains and
  Strings (1976), pp.~1387--1390.}
  {kibble-topology-defects-1976-ch23-v2ch23-p01}
  \begin{enumerate}
    \item[(a)]
    In a nonlinear sigma model whose field is constrained to a vacuum
    manifold \(\mathcal M\), let a localized texture approach the same chosen
    point of \(\mathcal M\) in every direction at spatial infinity in
    \(d\) dimensions.  Explain why a nonsingular configuration with
    this fixed boundary condition defines a based map
    \(S^d\to\mathcal M\).  Show that continuous deformations within this
    configuration space preserve its element of \(\pi_d(\mathcal M)\),
    and state precisely what must happen for this topological label to
    change.
    \item[(b)]
    A static defect has codimension \(q\), so a small surface linking it is
    \(S^{q-1}\).  Derive the rule that defects are classified by
    \(\pi_{q-1}(\mathcal M)\).  Apply it in three spatial dimensions to
    domain walls, strings, and monopoles.  For
    \(\mathcal M=S^1\), \(S^2\), and two disconnected points, identify
    which defect is topologically allowed in each case.
    \item[(c)]
    For a real scalar in one spatial dimension with
    \[
      E=\int dx\left[
         \frac12(\partial_x\phi)^2+
         \frac\lambda4(\phi^2-v^2)^2\right],
    \]
    define a topological charge using the values of \(\phi\) at
    \(x=\pm\infty\).  Complete the square to obtain a lower bound on the
    energy, solve the first-order equation for a wall centered at \(x_0\),
    and calculate its tension.
    \item[(d)]
    A complex scalar with vacuum manifold \(S^1\) has, far from a straight
    string, \(\Phi=v e^{in\varphi}\).  Prove that \(n\in\mathbb Z\) is
    invariant under continuous finite-energy deformations.  Using the
    normalization
    \(E/L=\int d^2x\,\lvert\nabla\Phi\rvert^2\), estimate the
    large-distance gradient energy per unit length between a core radius
    \(r_c\) and an infrared cutoff \(L\), and explain the logarithm.
    \item[(e)]
    Suppose a symmetry-breaking transition has correlation length
    \(\xi\), so regions separated by more than \(\xi\) choose their vacuum
    values approximately independently.  Explain how interpolation among
    several domains can force a noncontractible map into the vacuum
    manifold and thereby trap a defect.  In three spatial dimensions,
    estimate the wall area, string length, and monopole number per unit
    volume in terms of \(\xi\), up to dimensionless formation
    probabilities.
  \end{enumerate}

  \SupplementaryExercise{2}{Derrick Scaling with Several Derivative Orders}
  {Inspired by G. H. Derrick, Comments on Nonlinear Wave Equations as
  Models for Elementary Particles (1964), pp.~1252--1254.}
  {derrick-nonlinear-wave-1964-ch23-v2ch23-p02}
  In \(d\) Euclidean dimensions suppose a positive action is a sum
  \(I(R)=\sum_k R^{d-k}I_k\), where \(I_k\) contains \(k\) derivatives
  and \(I_k>0\).  Derive the stationarity and stability conditions under
  a scale transformation.  Show that a finite stable size requires
  nonzero terms on both sides of \(k=d\), and find the stationary radius
  when only terms of orders \(k_-<d<k_+\) are present.

  \SupplementaryExercise{3}{Monopole Topology and Gauge-Invariant Magnetic Charge}
  {Inspired by G. 't Hooft, Magnetic Monopoles in Unified Gauge
  Theories (1974), pp.~276--279.}
  {t-hooft-magnetic-monopoles-1974-ch23-v2ch23-p03}
  \begin{enumerate}
    \item[(a)]
    In an \(SU(2)\) gauge theory with an adjoint scalar
    \(\phi^a\), assume \(\phi^a\phi^a\to v^2\) at spatial infinity.
    Show that \(\widehat\phi^a=\phi^a/\lvert\phi\rvert\) defines a map
    \(S^2_\infty\to S^2\).  Write its integer degree as a surface
    integral and explain why a smooth finite-energy deformation cannot
    change it.
    \item[(b)]
    Define
    \[
      \mathcal F_{\mu\nu}
        =\widehat\phi^aF_{\mu\nu}^a
         -\frac1e\epsilon^{abc}\widehat\phi^a
          (D_\mu\widehat\phi)^b(D_\nu\widehat\phi)^c .
    \]
    Prove that \(\mathcal F_{\mu\nu}\) is gauge invariant.  In a patch
    where \(\widehat\phi^a=\delta^{a3}\), reduce it to the Abelian field
    strength of \(A_\mu^3\).
    \item[(c)]
    For a finite-energy monopole, use \(D_i\widehat\phi\to0\) at infinity
    to express the magnetic flux of \(\mathcal F_{ij}\) in terms of
    \(\widehat\phi\).  Show that
    \[
      g_m=\int_{S^2_\infty}\mathbf B\cdot d\mathbf S
         =\frac{4\pi n}{e},
    \]
    with \(n\) the degree from part~(a).  Relate this to Dirac
    quantization for particles whose smallest electric charge is \(e/2\).
  \end{enumerate}

  \SupplementaryExercise{4}{The Bogomolny Bound and Monopole Scales}
  {Inspired by E. B. Bogomolny, The Stability of Classical Solutions
  (1976), pp.~449--454.}
  {bogomolny-stability-classical-solutions-1976-ch23-v2ch23-p04}
  \begin{enumerate}
    \item[(a)]
    In the vanishing-scalar-potential limit, take
    \[
      E=\frac12\int d^3x\,
         \left(B_i^aB_i^a+D_i\phi^aD_i\phi^a\right).
    \]
    Complete the square to prove
    \(E\ge 4\pi v\lvert n\rvert/e\).  State the first-order equation that
    saturates the bound and identify the sign appropriate to \(n>0\).
    \item[(b)]
    In an adjoint-Higgs theory with gauge coupling \(e\), vacuum value
    \(v\), and scalar self-coupling \(\lambda\), use dimensional analysis
    and the particle masses \(m_W=ev\) and
    \(m_H=\sqrt{2\lambda}\,v\) to show that a unit monopole mass has the
    form
    \[
      M_{\mathrm{mon}}=\frac{4\pi v}{e}
        f\left(\frac{\lambda}{e^2}\right).
    \]
    Determine \(f(0)\) from the Bogomolny bound and identify the two
    length scales controlling the gauge and scalar tails.
  \end{enumerate}

  \SupplementaryExercise{5}{Regular BPS Monopole Profiles}
  {Inspired by M. K. Prasad and C. M. Sommerfield, Exact Classical
  Solution for the 't Hooft Monopole and the Julia--Zee Dyon (1975),
  pp.~760--762.}
  {prasad-sommerfield-monopole-1975-ch23-v2ch23-p05}
  For \(x=evr\), consider
  \[
    \phi^a=v\widehat r^a h(x),\qquad
    A_i^a=\epsilon_{aij}\widehat r^j
      \frac{1-K(x)}{er},
    \qquad
    K(x)=\frac{x}{\sinh x},\quad
    h(x)=\coth x-\frac1x .
  \]
  Find the small- and large-\(x\) behavior of \(K\) and \(h\).
  Demonstrate that the fields are regular at the origin, approach the
  broken-symmetry vacuum at infinity, and have a core radius of order
  \((ev)^{-1}\).

  \SupplementaryExercise{6}{Instanton Topology, Self-Duality, and the BPST Modulus}
  {Inspired by A. A. Belavin, A. M. Polyakov, A. S. Schwartz, and
  Yu. S. Tyupkin, Pseudoparticle Solutions of the Yang--Mills
  Equations (1975), pp.~85--87.}
  {bpst-pseudoparticle-yang-mills-1975-ch23-v2ch23-p06}
  \begin{enumerate}
    \item[(a)]
    With
    \(\widetilde F_{\mu\nu}
      =\tfrac12\epsilon_{\mu\nu\rho\sigma}F_{\rho\sigma}\),
    define
    \[
      Q=\frac1{16\pi^2}\int d^4x\,
         \operatorname{tr}(F_{\mu\nu}\widetilde F_{\mu\nu}),
      \qquad
      \operatorname{tr}(T^aT^b)=\frac12\delta^{ab}.
    \]
    Show that \(\operatorname{tr}(F\wedge F)\) is an exterior derivative.
    Explain why a finite-action gauge field that is pure gauge at infinity
    makes \(Q\) the winding number of a map \(S^3_\infty\to SU(2)\).
    \item[(b)]
    For Euclidean Yang--Mills theory,
    \(S_E=(2g^2)^{-1}\int d^4x\,\operatorname{tr}F_{\mu\nu}F_{\mu\nu}\),
    complete the square using \(F\mp\widetilde F\).  Prove
    \[
      S_E\ge\frac{8\pi^2}{g^2}\lvert Q\rvert
    \]
    and show that equality requires a self-dual or anti-self-dual field.
    Verify directly that such a field satisfies the Yang--Mills equation.
    \item[(c)]
    An \(SU(2)\) unit instanton centered at the origin has
    \[
      F_{\mu\nu}^a
        =-\frac{4\rho^2\eta^a_{\mu\nu}}
         {(x^2+\rho^2)^2},
      \qquad
      \eta^a_{\mu\nu}
        =\frac12\epsilon_{\mu\nu\rho\sigma}
          \eta^a_{\rho\sigma},
    \]
    where
    \(\eta^a_{\mu\nu}\eta^a_{\mu\nu}=12\).
    Compute \(F_{\mu\nu}^aF_{\mu\nu}^a\), integrate it over
    \(\mathbb R^4\), and recover \(Q=1\) and
    \(S_E=8\pi^2/g^2\).
    \item[(d)]
    Use the action density in part~(c) to show that changing
    \(\rho\) redistributes the instanton density without changing either
    \(Q\) or \(S_E\).  Calculate the radius \(R_{1/2}\) of a four-ball
    containing half of the action, and explain why classical scale
    invariance makes \(\rho\) a collective coordinate rather than a
    fixed parameter.
  \end{enumerate}

  \SupplementaryExercise{7}{Small-Instanton Convergence}
  {Inspired by Claude W. Bernard, Gauge Zero Modes, Instanton
  Determinants, and Quantum-Chromodynamic Calculations (1979),
  one-instanton measure.}
  {bernard-instanton-zero-modes-1979-ch23-v2ch23-p07}
  For pure \(SU(N)\), \(N\ge2\), suppose the one-instanton size measure
  has the one-loop form
  \[
    dn\propto d\rho\,\rho^{-5}(\rho\mu)^{b_0}
      \exp\left[-\frac{8\pi^2}{g^2(\mu)}\right],
    \qquad b_0=\frac{11N}{3}.
  \]
  Find the small-\(\rho\) power of the integrand and determine whether
  the size integral converges at \(\rho=0\).  Analyze its behavior at
  large \(\rho\) and explain why that region signals a breakdown of the
  dilute semiclassical calculation rather than a physical divergence
  that can be trusted quantitatively.

  \SupplementaryExercise{8}{The Dilute Instanton Gas and Vacuum Angle}
  {Inspired by C. G. Callan Jr., R. F. Dashen, and D. J. Gross,
  The Structure of the Gauge Theory Vacuum (1976), pp.~334--340.}
  {callan-dashen-gross-gauge-vacuum-1976-ch23-v2ch23-p08}
  Suppose each instanton or anti-instanton has spacetime-volume weight
  \(K e^{-8\pi^2/g^2}\) and phases \(e^{\pm i\theta}\).  Sum over
  independent numbers \(n_+\) and \(n_-\) to obtain the dilute-gas
  partition function.  Extract the \(\theta\)-dependent vacuum-energy
  density and the leading topological susceptibility.

  \SupplementaryExercise{9}{Strong \(CP\), Topological Susceptibility, and Axion Observables}
  {Inspired by the axion mechanism in Steven Weinberg, A New Light
  Boson? (1978), pp.~223--226; extended editorial formulation.}
  {weinberg-new-light-boson-1978-ch23-v2ch23-p09}
  \begin{enumerate}
    \item[(a)]
    In QCD with quark mass matrix \(M_q\), perform independent anomalous
    axial rephasings of the quarks.  Show that the phase of
    \(\det M_q\) and the coefficient \(\theta\) of
    \(F_{\mu\nu}^a\widetilde F^{a\mu\nu}\) shift oppositely.  Identify the
    invariant angle that can enter observables.
    \item[(b)]
    Suppose \(m_u=0\) exactly while all other quarks are massive.  Use an
    anomalous axial rotation of \(u\) to remove the strong vacuum angle.
    Explain why this does not introduce a compensating phase in the mass
    matrix, and show how the result appears in the light-quark expression
    for the topological susceptibility.
    \item[(c)]
    For two light quarks use
    \[
      V=-v(m_u\cos\alpha_u+m_d\cos\alpha_d),\qquad
      \alpha_u+\alpha_d=\theta.
    \]
    Minimize to quadratic order in \(\theta\) and derive
    \(\chi=\partial_\theta^2E(\theta)|_{\theta=0}\).
    Express the result in terms of \(m_\pi\) and \(F_\pi\), and check the
    equal-mass and one-massless-quark limits.
    \item[(d)]
    For
    \[
      \mathcal L_{\mathrm{int}}
        =-\frac{g_{a\gamma\gamma}}4
         aF_{\mu\nu}\widetilde F^{\mu\nu},
    \]
    derive the decay rate \(\Gamma(a\to\gamma\gamma)\).  Then set
    \(g_{a\gamma\gamma}=C_\gamma\alpha/(2\pi f_a)\) and use
    \(m_af_a=\Lambda_a^2\) to express the rate entirely in terms of
    \(m_a\), \(\Lambda_a\), \(C_\gamma\), and \(\alpha\).  Determine the
    mass scaling of the lifetime.
    \item[(e)]
    Let
    \[
      V=-v\left[
        m_u\cos\left(\frac{\pi^0}{F_\pi}-c_u\frac a{f_a}\right)
        +m_d\cos\left(\frac{\pi^0}{F_\pi}+c_d\frac a{f_a}\right)
        \right],
      \qquad c_u+c_d=1.
    \]
    Expand to quadratic order and find the mass matrix.  In the limit
    \(f_a\gg F_\pi\), calculate the axion mass and the leading pion
    component of the light eigenstate.  Show that the mass is independent
    of the arbitrary split between \(c_u\) and \(c_d\).
  \end{enumerate}

  \SupplementaryExercise{10}{Axion Periodicity and Domain Walls}
  {Inspired by Pierre Sikivie, Axions, Domain Walls, and the Early
  Universe (1982), pp.~1156--1159.}
  {sikivie-axion-domain-walls-1982-ch23-v2ch23-p10}
  Suppose the QCD anomaly of the Peccei--Quinn current is the integer
  \(N_{\mathrm{DW}}\), so the low-energy potential has the form
  \[
    V(a)=\chi\left[
      1-\cos\left(N_{\mathrm{DW}}\frac a{f_a}+\theta\right)\right],
    \qquad a\sim a+2\pi f_a.
  \]
  Locate all physically distinct minima, find the small-oscillation
  mass, and explain why \(N_{\mathrm{DW}}>1\) permits stable domain
  walls whereas \(N_{\mathrm{DW}}=1\) does not leave distinct vacua.

  \SupplementaryExercise{11}{Euclidean Bounces and Lorentzian Bubble Evolution}
  {Inspired by Sidney Coleman, Fate of the False Vacuum:
  Semiclassical Theory (1977), pp.~2929--2932.}
  {coleman-false-vacuum-1977-ch23-v2ch23-p11}
  \begin{enumerate}
    \item[(a)]
    For a scalar Euclidean action
    \[
      S_E=\int d^4x\left[\frac12(\partial_\mu\phi)^2+V(\phi)\right],
    \]
    assume the relevant bounce is \(O(4)\)-symmetric,
    \(\phi=\phi(\rho)\).  Derive its differential equation and boundary
    conditions.  Interpret the equation as one-dimensional motion in
    \(-V\) with a friction term, and explain the overshoot/undershoot
    criterion for finding the bounce.
    \item[(b)]
    Subtract the false-vacuum action and write the bounce exponent as
    \(B=T+U\), where
    \[
      T=\frac12\int d^4x\,(\partial\phi)^2,\qquad
      U=\int d^4x\,[V(\phi)-V(\phi_f)].
    \]
    Scale \(\phi(x)\to\phi(x/R)\) and use stationarity at \(R=1\) to
    derive a virial identity.  Express \(B\) in terms of either \(T\) or
    \(U\), prove that \(U<0\), and interpret this integrated statement.
    \item[(c)]
    Let an \(O(4)\)-symmetric Euclidean bounce be
    \(\phi_b(\rho)\), with
    \(\rho=\sqrt{r^2+\tau^2}\).  Continue \(\tau\to it\) and show that the
    Lorentzian field has initial data
    \(\phi(r,0)=\phi_b(r)\) and \(\dot\phi(r,0)=0\).  Explain how the
    Euclidean \(O(4)\) symmetry becomes Lorentzian \(O(3,1)\) symmetry and
    why the critical bubble is said to materialize at rest.
    \item[(d)]
    In the thin-wall approximation, suppose the Euclidean wall lies at
    \(\rho=R\).  After \(\tau\to it\), derive
    \(r^2-t^2=R^2\).  Find \(r(t)\), the wall velocity, its Lorentz factor,
    and its proper acceleration.  Show that the expanding wall becomes
    asymptotically lightlike.
    \item[(e)]
    A spherical thin wall of radius \(r\), tension \(\sigma\), and Lorentz
    factor \(\gamma\) has energy
    \(E_{\mathrm{wall}}=4\pi r^2\sigma\gamma\).  The true-vacuum interior
    releases
    \(E_{\mathrm{vol}}=(4\pi/3)\epsilon r^3\).  Use
    \(\gamma=r/R\) and the nucleation radius
    \(R=3\sigma/\epsilon\) to show that these are equal.  Interpret the
    vanishing total energy relative to the false vacuum.
  \end{enumerate}

  \SupplementaryExercise{12}{The Unique Negative Bounce Mode}
  {Inspired by Curtis G. Callan Jr. and Sidney Coleman, Fate of the
  False Vacuum. II. First Quantum Corrections (1977), pp.~1762--1768.}
  {callan-coleman-false-vacuum-ii-1977-ch23-v2ch23-p12}
  For the thin-wall function
  \[
    B(R)=2\pi^2\sigma R^3-\frac{\pi^2}{2}\epsilon R^4,
  \]
  evaluate its second derivative at the critical radius and identify
  the associated fluctuation.  Explain why one negative eigenvalue
  produces an imaginary part of the false-vacuum energy, why four
  translational zero modes instead produce a spacetime-volume factor,
  and how the sign of the imaginary part is fixed to give a positive
  decay rate.

  \SupplementaryExercise{13}{Anomaly-Induced Domain-Wall Charge}
  {Adapted from John McGreevy, Physics 215C, Assignment 2, Problem 4 (2025).}
  {mcgreevy-215c-s25-a2-p4-v2ch23-p13}
  A Dirac fermion in \(1+1\) dimensions has
  \[
    \mathcal L
      =\overline\psi\left(i\slashed\partial
        -m e^{i\gamma_5\theta(x)}\right)\psi ,
    \qquad j^\mu=q\overline\psi\gamma^\mu\psi .
  \]
  \begin{enumerate}
    \item[(a)] Explain why the vacuum current vanishes when \(\theta\) is
      constant.
    \item[(b)] Couple \(j^\mu\) to a nondynamical \(A_\mu\), define
      \(\Gamma[A,\theta]\) by integrating out the fermion, and show that
      the term linear in \(A_\mu\) is the vacuum current.
    \item[(c)] Perform a local chiral rotation that makes the mass real.
      Track both the derivative coupling and the anomalous Jacobian.
    \item[(d)] To leading order in derivatives, derive
      \(\langle j^\mu\rangle\) in terms of \(\partial_\nu\theta\), and show
      why the coefficient is independent of \(m\).
    \item[(e)] Compute the fermion number bound to a static wall with
      \(\theta(-\infty)=0\) and \(\theta(+\infty)=\pi\).  Explain the
      integer ambiguity and the meaning of a \(2\pi\) interpolation.
    \item[(f)] For a real mass profile \(M(x)\) that changes sign across
      the wall, solve the associated Dirac equation at zero energy and
      exhibit the exponentially localized mode.
  \end{enumerate}

  \SupplementaryExercise{14}{Peierls Dimerization and the Polyacetylene Wall Mode}
  {Adapted from John McGreevy, Physics 215C, Assignment 3, Problem 1 (2025).}
  {mcgreevy-215c-s25-a3-p1-v2ch23-p14}
  Ignore electron spin and consider an even periodic chain with slow ionic
  displacements \(u_n\):
  \[
    H=-t\sum_n\left[(1+u_n)c_n^\dagger c_{n+1}
      +\mathrm{H.c.}\right]
      +K\sum_n(u_n-u_{n+1})^2 .
  \]
  \begin{enumerate}
    \item[(a)] For \(u_n=\phi(-1)^n\), use a two-site unit cell to obtain
      both electronic bands.  Show that nonzero \(\phi\) opens a gap, and
      expand near a band-touching point to identify the continuum Dirac
      velocity and mass.
    \item[(b)] At half filling, compute the electronic ground-state energy
      as a function of \(\phi\).  Check the uniform-chain and isolated-dimer
      limits.  Add the elastic energy, analyze the total energy for small
      \(\phi\), and find the nonzero minimizing dimerization in the regime
      where that expansion is valid.
    \item[(c)] Explain why the two signs of \(\phi\) are degenerate.
      Construct a wall between them and show that it carries a midgap
      fermion whose empty and filled states have charges
      \(-q/2\) and \(+q/2\).
    \item[(d)] Verify the wall state directly in the lattice hopping
      matrix, giving the zero-energy recursion and its localization
      criterion.  Explain the finite-size behavior for periodic boundary
      conditions.
    \item[(e)] Allow complex hopping amplitudes.  Explain in both the
      continuum and lattice descriptions how the mass can then interpolate
      between the two real dimerizations without a protected midgap state.
  \end{enumerate}

  \SupplementaryExercise{15}{The Complete \(\phi^4\) Kink}
  {Adapted from Cambridge Part III, Solitons Example Sheet 1, Question 1.}
  {cambridge-part-iii-solitons-3p11a-q1-v2ch23-p15}
  For a real scalar in \(1+1\) dimensions with
  \[
    V(\phi)=\frac{\lambda}{4}(\phi^2-v^2)^2,
  \]
  derive the mechanical first integral for a finite-energy static field,
  rewrite its energy as a Bogomolny square plus a boundary term, and solve
  for the kink from \(-v\) to \(+v\).  Compute its mass and leading
  exponential approach to each vacuum.  Show explicitly that the
  first-order Bogomolny equation implies the second-order field equation,
  and interpret the latter as one-dimensional mechanical motion in the
  inverted potential \(-V\).  Finally, explain why a finite-energy
  first-order solution cannot skip an intermediate vacuum of a potential
  with three degenerate minima.

  \SupplementaryExercise{16}{Gauge-Invariant Vortex Winding and Radial Gauge}
  {Adapted from Cambridge Part III, Solitons Example Sheet 1, Question 8.}
  {cambridge-part-iii-solitons-3p11a-q8-v2ch23-p16}
  In the Abelian Higgs model on the plane, let a finite-energy configuration
  obey
  \[
    \frac{\phi}{|\phi|}\longrightarrow e^{in\vartheta}
    \quad (r\longrightarrow\infty),\qquad
    D_i=\partial_i-iqA_i .
  \]
  \begin{enumerate}
    \item[(a)] Show that a smooth, single-valued gauge transformation on the
      whole plane has zero net winding on a sufficiently large circle.
      Deduce that the integer \(n\) is gauge invariant.
    \item[(b)] Under
      \(\phi\mapsto e^{i\alpha}\phi\) and
      \(A_i\mapsto A_i+q^{-1}\partial_i\alpha\), construct a smooth gauge
      transformation that sets \(A_r=0\).  Explain why one cannot in
      general set \(A_\vartheta=0\) as well for a vortex of nonzero flux.
  \end{enumerate}

  \SupplementaryExercise{17}{The Critical Abelian-Higgs Bogomolny Bound}
  {Adapted from Cambridge Part III, Solitons Example Sheet 1, Question 9.}
  {cambridge-part-iii-solitons-3p11a-q9-v2ch23-p17}
  First verify
  \[
    \partial_i(\phi^*D_j\phi)
      =(D_i\phi)^*D_j\phi+\phi^*D_iD_j\phi,
    \qquad [D_i,D_j]\phi=-iqF_{ij}\phi .
  \]
  Then, at \(\lambda=q^2\), complete the Abelian-Higgs energy into squares.
  Derive the first-order vortex equations and prove
  \[
    E\geq2\pi v^2|n|.
  \]
  Show that the first-order equations imply the second-order field
  equations, and explain why well-separated critical vortices exert no
  static force.

\end{enumerate}
